\chapter{Requirements}
\label{section:Requirements}

We differentiate between functional and non-functional requirements \autocite{Herrmann2022}.

\section{Functional Requirements}
\label{subsection:FunctionalRequirements}

\begin{enumerate}
	\item[interoperability] The memory format is standardized.
	\item[columnar] A columnar format has advantages outlined in \autoref{section:column_vs_row}.
	\item[feature toggle] We must be able to (de)activate the new features at runtime.
	      This makes comparing the new features easier.
	\item[compatibility]  The electric vehicles example must complete without errors.
	      This example is executed, by running the \ac{CLI} command \Verb|npm run example:vehicles-polars| inside the project directory.
	      The \Verb|.sqlite| files produced by the different implementations must be identical.
	\item[modularization] Should it be necessary to write code in a language other than \emph{TypeScript}, this code must be placed in an external library usable from \emph{TypeScript}.
	\item[extensibility] The functionality of the chosen data representation must be extensible.
\end{enumerate}

\section{Non-Functional Requirements}
\label{subsection:NonFunctionalRequirements}

\begin{enumerate}
	\item[performance] The main goal of this thesis is to optimize \emph{Jayvee}'s in-memory representation.
	      We expect this to result in either a decrease in the execution time for \emph{Jayvee} models, or an extension of the \emph{Jayvee} interpreter's capabilities.
	\item[code style] The source code adheres to the project's code style.
	      This is enforced by running \Verb|npm run lint|.
	      In rust code we use the defaults from \emph{rustfmt} and \emph{clippy}.
	\item[maturity] We aim to create a prototype, not a mature implementation.
\end{enumerate}


Both types of requirements have been evaluated for completeness, consistency and feasibility by us.
